% ============================================================
% REGISTRO DE EXÁMENES
% Procedimiento funcional basado en RegistrarExamen.jsx
% Evidencias: cuatro etapas, escritorio y móvil
% ============================================================

\subsection{Registrar un examen: procedimiento completo}
\label{sec:registro-examen}

El registro de un examen se realiza desde
\textbf{Exámenes} mediante la opción
\textbf{Registrar examen}. La aplicación organiza el
procedimiento en cinco etapas: \textbf{Examen},
\textbf{Grupos}, \textbf{Aulas},
\textbf{Quién atiende} y \textbf{Listo}.

Las primeras cuatro etapas permiten preparar los datos.
La última muestra el resultado después de guardar
correctamente el examen.

\medskip
\begin{center}
\small\bfseries\color{Teal}
01 Examen
\quad--\quad
02 Grupos
\quad--\quad
03 Aulas
\quad--\quad
04 Personal
\quad--\quad
05 Listo
\end{center}
\medskip

\noindent\textbf{Antes de empezar.}
La cuenta debe tener permiso para gestionar exámenes,
disponer de una asignatura habilitada y trabajar
dentro del período académico correspondiente.

\noindent\textbf{Importante.}
La disponibilidad de asignaturas, grupos, ambientes
y personal depende de la información académica
registrada en el sistema.

% ============================================================
% MACRO EDITORIAL PARA LOS PARES DE CAPTURAS
% ============================================================

\newcommand{\CapturasRegistro}[3]{%
    \par\medskip
    \noindent
    \begin{minipage}{\linewidth}
        \centering
        \includegraphics[
            width=0.98\linewidth,
            height=0.39\textheight,
            keepaspectratio
        ]{imagenes/docente/flujo-registro/escritorio/paso-#1.png}
        \captionof{figure}{#2: vista de escritorio.}
        \label{fig:registro-escritorio-#1}
    \end{minipage}

    \par\medskip

    \noindent
    \begin{minipage}[t]{0.42\linewidth}
        \vspace{0pt}
        \centering
        \includegraphics[
            width=\linewidth,
            height=0.40\textheight,
            keepaspectratio
        ]{imagenes/docente/flujo-registro/movil/paso-#1.png}
        \captionof{figure}{#2: vista móvil.}
        \label{fig:registro-movil-#1}
    \end{minipage}
    \hfill
    \begin{minipage}[t]{0.53\linewidth}
        \vspace{3pt}
        {\bfseries\color{Teal}Uso desde el teléfono}\par
        \small #3
    \end{minipage}

    \par\medskip
}

% ============================================================
% ETAPA 1
% ============================================================

\subsubsection{Etapa 1. Datos generales del examen}

La primera pantalla determina qué evaluación se
realizará, a qué asignatura pertenece y cuándo tendrá
lugar. Estos datos condicionan las opciones de las
etapas posteriores.

\noindent\textbf{Procedimiento}

\begin{enumerate}
    \item Seleccione la \textbf{asignatura}
    que corresponde al examen.

    \item Elija el \textbf{tipo de evaluación}:
    primer parcial, segundo parcial, final,
    segunda instancia o mesa, según corresponda.

    \item Indique la \textbf{fecha}.
    Compruebe que pertenece al período y al
    calendario permitido para ese tipo de examen.

    \item Establezca la \textbf{hora de inicio}
    y la \textbf{duración} en minutos.

    \item Revise las \textbf{normas del examen}.
    Si se han recuperado normas de una evaluación
    anterior de la misma asignatura, compruebe que
    siguen siendo adecuadas.

    \item Pulse \textbf{Continuar}.
\end{enumerate}

\noindent\textbf{Validaciones que aplica el sistema}

\begin{itemize}
    \item La fecha y la hora son obligatorias.
    \item No se admite programar una nueva evaluación
    en una fecha anterior a la actual.
    \item La hora de inicio debe encontrarse
    entre las 06:00 y las 22:00.
    \item La duración debe estar entre
    15 y 480 minutos.
    \item El examen no debe terminar después
    de medianoche.
    \item Deben respetarse las fechas del período
    y la ventana académica correspondiente.
\end{itemize}

Si algún dato incumple estas restricciones,
el sistema informa el motivo y mantiene al usuario
en la etapa para que lo corrija.

\CapturasRegistro
{1}
{Etapa 1, información general y normas}
{Los campos se organizan verticalmente.
Desplácese para revisar las normas y localizar
el botón Continuar. No avance sin comprobar
fecha, hora y duración.}

% ============================================================
% ETAPA 2
% ============================================================

\subsubsection{Etapa 2. Selección de grupos}

Un examen debe vincularse a los grupos académicos
que participarán. Esta selección determina parte
del alcance de estudiantes y las opciones posteriores
de asignación de ambientes.

\noindent\textbf{Procedimiento}

\begin{enumerate}
    \item Compruebe que aparece la asignatura
    seleccionada en la primera etapa.

    \item Localice los grupos académicos disponibles.

    \item Seleccione los grupos que rendirán
    la evaluación.

    \item Si hay numerosos grupos, utilice
    el buscador disponible para encontrarlos.

    \item Revise la selección antes de continuar.
\end{enumerate}

\noindent\textbf{Comprobación}

El formulario debe conservar al menos un grupo
seleccionado. El código del grupo y la información
docente permiten comprobar que la selección
corresponde a los participantes previstos.

No seleccione grupos ajenos al examen únicamente
para avanzar en el formulario.

\CapturasRegistro
{2}
{Etapa 2, selección de grupos académicos}
{Revise las casillas de selección.
En pantallas estrechas puede ser necesario
desplazarse para consultar el listado completo
y las acciones de navegación.}

% ============================================================
% ETAPA 3
% ============================================================

\subsubsection{Etapa 3. Asignación de aulas}

En esta etapa se eligen los ambientes en los
que se realizará la evaluación. Las opciones
disponibles se consultan tomando en cuenta
los grupos seleccionados.

\noindent\textbf{Procedimiento}

\begin{enumerate}
    \item Abra la etapa \textbf{Aulas}.

    \item Examine los ambientes que presenta
    el sistema.

    \item Identifique cada aula por su nombre
    y la ubicación informada.

    \item Seleccione los ambientes necesarios
    para organizar el examen.

    \item Compruebe que la selección coincide
    con la planificación de ambientes.

    \item Continúe a la asignación de personal.
\end{enumerate}

\noindent\textbf{Comprobación}

Debe existir al menos un ambiente seleccionado
para continuar con un examen que requiere control
de ingreso por aula.

\noindent\textbf{Precaución}

La selección visual no sustituye la revisión de
la planificación académica ni la comprobación
de restricciones que informe el servidor.

\CapturasRegistro
{3}
{Etapa 3, selección de ambientes}
{La vista reorganiza los controles de aulas.
Compruebe los nombres de los ambientes seleccionados;
si existe un panel lateral, desplácese para consultarlo.}

% ============================================================
% ETAPA 4
% ============================================================

\subsubsection{Etapa 4. Personal responsable}

Cada aula seleccionada debe relacionarse con
la información del personal que atenderá
el examen. El formulario contempla una persona
encargada y personal de apoyo cuando corresponda.

\noindent\textbf{Procedimiento}

\begin{enumerate}
    \item Acceda a \textbf{Quién atiende}.

    \item Identifique las aulas que fueron
    seleccionadas previamente.

    \item Revise la persona encargada de
    cada ambiente.

    \item Compruebe la asignación del personal
    de apoyo correspondiente.

    \item Verifique que no existan asignaciones
    incorrectas antes de registrar.

    \item Pulse \textbf{Registrar examen}
    si se trata de uno nuevo, o
    \textbf{Guardar examen} si está editando
    un registro propio existente.
\end{enumerate}

\noindent\textbf{Importante}

Al guardar, la aplicación envía los datos al
servidor. Si este rechaza algún campo, el
formulario vuelve a la etapa donde debe
corregirse el problema.

Esta operación modifica información del sistema;
no debe ejecutarse únicamente para obtener
capturas de documentación.

\CapturasRegistro
{4}
{Etapa 4, encargados y personal de apoyo}
{Revise la correspondencia entre cada aula
y sus responsables. La acción de guardar
debe utilizarse solamente cuando la información
haya sido verificada.}

% ============================================================
% ETAPA 5
% ============================================================

\subsubsection{Etapa 5. Confirmación y resumen}

Cuando el servidor registra o actualiza
correctamente el examen, la aplicación muestra
la etapa \textbf{Listo}, con el resumen de
los datos guardados.

El usuario debe revisar:

\begin{enumerate}
    \item La asignatura y el tipo de evaluación.
    \item La fecha, hora y duración.
    \item Los grupos vinculados.
    \item Los ambientes asignados.
    \item El personal responsable de cada aula.
    \item Las normas configuradas para el examen.
\end{enumerate}

\noindent\textbf{Resultado esperado}

El sistema confirma el registro y permite
consultar la información guardada.

\noindent\textbf{Si aparece un error}

No considere completada la operación
hasta que el sistema confirme el guardado.
Regrese al campo señalado, corrija el dato
y vuelva a comprobar la respuesta.

\medskip
\noindent
{\color{Teal}\bfseries Continuación del proceso.}
Una vez creado el examen, la preparación
puede continuar con la habilitación de
estudiantes y los códigos QR, siempre que
la cuenta disponga de los permisos necesarios.

\begin{notebox}
\textbf{Cobertura visual de este procedimiento.}

Las figuras anteriores corresponden a
las cuatro etapas editables de un examen
existente, capturadas sin guardar cambios.
La confirmación de una operación nueva
debe verificarse funcionalmente antes de
añadir su captura como evidencia de registro.
\end{notebox}

% FIN DEL CAPÍTULO DE REGISTRO
